% ============================================================
\chapter{Liens avec la TDA}
\label{ch:tda}
% ============================================================

L'apprentissage g\'eom\'etrique et l'analyse topologique des donn\'ees (TDA) sont deux r\'eponses compl\'ementaires \`a la m\^eme question~: comment exploiter la \emph{structure} des donn\'ees~? Les GNN capturent les relations locales et les sym\'etries~; la TDA extrait des invariants topologiques globaux (nombres de Betti, homologie persistante) qui sont robustes au bruit et invariants par d\'eformation continue. Combiner les deux --- par exemple en ajoutant des features topologiques aux repr\'esentations d'un GNN, ou en utilisant des couches de persistance diff\'erentiables --- am\'eliore l'expressivit\'e et la performance. Ce chapitre explore ces connexions, de la th\'eorie (limites d'expressivit\'e des GNN r\'ev\'el\'ees par la topologie) aux applications (filtration apprise, couches de persistance).

\begin{intuition}
L'analyse topologique des donn\'ees (TDA)\index{TDA} et
l'apprentissage g\'eom\'etrique partagent une pr\'eoccupation
commune~: exploiter la \textbf{structure} des donn\'ees. La TDA
apporte des invariants topologiques (nombres de Betti, homologie
persistante) qui compl\`etent les features apprises par les GNN
et am\'eliorent leur \textbf{expressivit\'e}.
\end{intuition}

% ============================================================
\section{Rappels d'homologie persistante}
\label{sec:rappels_persistance}
% ============================================================

\begin{definition}[Homologie persistante]
\index{homologie persistante}
\'Etant donn\'ee une filtration de complexes simpliciaux
$K_0 \subseteq K_1 \subseteq \cdots \subseteq K_m$, l'\textbf{homologie
persistante} suit l'\'evolution des classes d'homologie \`a travers
la filtration. Chaque classe na\^\i t \`a un indice $b$ et meurt
\`a un indice $d > b$, donnant un point $(b, d)$ dans le
\textbf{diagramme de persistance}.
\end{definition}

\begin{remarque}
Pour un graphe $G = (V, E)$, on peut construire diff\'erentes
filtrations~:
\begin{itemize}
  \item \textbf{Par degr\'e}~: $f(v) = \mathrm{deg}(v)$ et
        $f(e) = \max(f(u), f(v))$.
  \item \textbf{Par centralit\'e}~: $f(v) = \mathrm{betweenness}(v)$.
  \item \textbf{Par features apprises}~: $f(v) = g_\theta(h_v)$
        o\`u $g_\theta$ est un r\'eseau de neurones.
\end{itemize}
\end{remarque}

% ============================================================
\section{Features topologiques dans les GNN}
\label{sec:features_topo}
% ============================================================

\begin{definition}[Features topologiques pour les graphes]
\index{features topologiques}
Les \textbf{features topologiques} d'un graphe sont extraites
de ses diagrammes de persistance~:
\begin{itemize}
  \item $H_0$~: les composantes connexes et leur stabilit\'e.
  \item $H_1$~: les cycles et leur persistance.
\end{itemize}
Ces features sont vectoris\'ees (paysages, images de persistance)
et concat\'en\'ees aux features du GNN.
\end{definition}

\begin{proposition}[Compl\'ementarit\'e]
Les features topologiques capturent des propri\'et\'es globales
du graphe (cycles, cavit\'es) que le message passing local
ne peut pas d\'etecter en un nombre born\'e d'it\'erations~:
d\'etecter un cycle de longueur $k$ n\'ecessite au moins $k/2$
couches de message passing.
\end{proposition}

\begin{minted}{python}
import torch
import numpy as np
from gudhi import RipsComplex
from gudhi.representations import PersistenceImage

def topo_features(pos, max_edge=5.0, resolution=10):
    """Calcule les features topologiques d'un nuage de points."""
    rips = RipsComplex(points=pos, max_edge_length=max_edge)
    st = rips.create_simplex_tree(max_dimension=2)
    st.persistence()

    # Diagrammes par dimension
    dgm_h0 = st.persistence_intervals_in_dimension(0)
    dgm_h1 = st.persistence_intervals_in_dimension(1)

    # Vectorisation par images de persistance
    pi = PersistenceImage(
        bandwidth=0.1, resolution=[resolution, resolution]
    )
    feats = []
    for dgm in [dgm_h0, dgm_h1]:
        if len(dgm) > 0:
            dgm = dgm[np.isfinite(dgm[:, 1])]
        if len(dgm) > 0:
            feats.append(pi.fit_transform([dgm])[0])
        else:
            feats.append(np.zeros(resolution ** 2))
    return np.concatenate(feats)
\end{minted}

% ============================================================
\section{Message passing topologique}
\label{sec:message_passing_topo}
% ============================================================

\begin{definition}[Message passing sur complexes simpliciaux]
\index{message passing simplicial}
Le \textbf{message passing simplicial} (Bodnar et al.\ 2021)
g\'en\'eralise le message passing des graphes aux complexes
simpliciaux de dimension sup\'erieure~:
\[
  h_\sigma^{(t+1)} = \phi\!\left(h_\sigma^{(t)},\;
  \bigoplus_{\tau \in \mathcal{B}(\sigma)} \psi_{\mathcal{B}}(h_\tau^{(t)}),\;
  \bigoplus_{\tau \in \mathcal{C}(\sigma)} \psi_{\mathcal{C}}(h_\tau^{(t)})
  \right)
\]
o\`u $\mathcal{B}(\sigma)$ sont les faces (boundary) et
$\mathcal{C}(\sigma)$ les co-faces (coboundary) du simplexe $\sigma$.
\end{definition}

\begin{intuition}
Dans un graphe, les messages circulent le long des ar\^etes
(1-simplexes). Dans un complexe simplicial, on ajoute des messages
le long des triangles (2-simplexes) et des structures de dimension
sup\'erieure. Cela permet de capturer des relations
\textbf{d'ordre sup\'erieur} entre les n\oe uds.
\end{intuition}

\begin{definition}[Message passing sur complexes cellulaires]
\index{CW Network}
Les \textbf{CW Networks} (Bodnar et al.\ 2021) \'etendent le
message passing aux complexes cellulaires, qui g\'en\'eralisent
les complexes simpliciaux en autorisant des cellules de forme
arbitraire (pas seulement des simplexes).
\end{definition}

% ============================================================
\section{Hi\'erarchie de Weisfeiler-Leman et topologie}
\label{sec:wl_topo}
% ============================================================

\begin{definition}[Test de Weisfeiler-Leman]
\index{Weisfeiler-Leman}
Le \textbf{test de Weisfeiler-Leman} (WL) de niveau $k$ est un
algorithme de raffinement de couleurs qui caract\'erise les graphes~:
deux graphes non distingu\'es par $k$-WL sont dits
\textbf{$k$-WL \'equivalents}.
\end{definition}

\begin{theoreme}[Expressivit\'e des GNN --- Xu et al.\ 2019, Morris et al.\ 2019]
\index{expressivit\'e WL}
Un GNN de type message passing est \textbf{au plus aussi puissant}
que le test $1$-WL pour distinguer des graphes non isomorphes.
\end{theoreme}

\begin{theoreme}[Au-del\`a de WL par la topologie --- Bodnar et al.\ 2021]
\index{complexe simplicial}
Le message passing sur des complexes simpliciaux avec des features
topologiques (homologie persistante) peut distinguer des graphes
que $k$-WL ne distingue pas, pour tout $k$ fix\'e.
\end{theoreme}

\begin{exemple}[Graphes de Rook]
Les graphes de Rook $R_{3 \times 3}$ et $R_{4 \times 4}$ ne sont
pas distingu\'es par $3$-WL. Cependant, leurs diagrammes de
persistance $H_1$ (avec filtration par degr\'e) diff\`erent~:
$R_{4 \times 4}$ a davantage de cycles persistants.
\end{exemple}

% ============================================================
\section{Expressivit\'e et topologie}
\label{sec:expressivite}
% ============================================================

\begin{proposition}[Enrichissement topologique]
En ajoutant des features d'homologie persistante aux features de
n\oe uds d'un GNN, l'expressivit\'e r\'esultante est
\textbf{strictement sup\'erieure} \`a celle du GNN seul pour
le probl\`eme de classification de graphes.
\end{proposition}

\begin{definition}[Filtration apprise jointe]
\index{filtration apprise}
Dans une \textbf{filtration apprise jointe}, les valeurs de
filtration sont d\'efinies par le GNN lui-m\^eme~:
\[
  f_\theta(v) = \mathrm{MLP}_\theta(h_v^{(L)})
\]
o\`u $h_v^{(L)}$ est l'embedding final du n\oe ud $v$ apr\`es $L$
couches de message passing. La persistance est calcul\'ee sur cette
filtration, et le gradient r\'etropropag\'e \`a travers le calcul
de persistance (cf.\ persistance diff\'erentiable).
\end{definition}

\begin{formules}[title=\textbf{Pipeline GNN + TDA}]
\[
  X \xrightarrow{\text{GNN}} H^{(L)}
  \xrightarrow{f_\theta} \text{Filtration}
  \xrightarrow{\text{PH}} \mathrm{Dgm}
  \xrightarrow{\text{PersLay}} z_{\mathrm{topo}}
  \xrightarrow{\oplus} \text{Pr\'ediction}
\]
\end{formules}

% ============================================================
\section{Exercices}
% ============================================================

\begin{exercice}
Pour le graphe de Petersen, calculer le diagramme de persistance
$H_0$ et $H_1$ avec la filtration par degr\'e. Tous les n\oe uds
ont le m\^eme degr\'e~: que peut-on en d\'eduire~?
\end{exercice}

\begin{exercice}
Montrer que deux graphes isomorphes ont les m\^emes diagrammes de
persistance pour toute filtration d\'efinie \`a partir de
propri\'et\'es structurelles du graphe.
\end{exercice}

\begin{exercice}[Message passing simplicial]
\'Ecrire les \'equations de message passing pour un triangle
$(v_1, v_2, v_3)$ dans un complexe simplicial. Quels messages
re\c{c}oit l'ar\^ete $(v_1, v_2)$~?
\end{exercice}

\begin{exercice}
Donner un exemple de deux graphes non isomorphes ayant les m\^emes
diagrammes de persistance pour la filtration par degr\'e. Proposer
une filtration qui les distingue.
\end{exercice}

\begin{exercice}[Projet]
Impl\'ementer un GNN enrichi de features topologiques sur le
benchmark ZINC et comparer avec un GCN de base. Mesurer
l'am\'elioration en MAE.
\end{exercice}
